% Source: physical PDF pp. 331--332 (printed pp. 308--309), from the
% Section 30.1 heading through the final sentence immediately before
% Section 30.2 on physical p. 333.
% Inventory: Eqs. (30.1.1)--(30.1.9), three unnumbered display groups,
% no citations, figures, tables, footnotes, or typographic dividers.
% Uncertainties: none.

\subsection{Potential Superfields}

Consider a theory of left-chiral superfields
\(\Phi_n(x,\theta,\bar\theta)\) and their complex conjugates, but for
simplicity without gauge superfields. All of the
vacuum expectation values of time-ordered products of component fields may be
calculated from vacuum expectation values of the time-ordered products of
these superfields. We might try to calculate these from the path-integral
formula
\begin{align}
  &\left\langle
    T\left\{\Phi_{n_1}(x_1,\theta_1,\bar\theta_1),
      \Phi_{n_2}(x_2,\theta_2,\bar\theta_2),\ldots\right\}
   \right\rangle \notag\\
  &\qquad={}
  \int\left[\prod_{n,x,\theta,\bar\theta}
      \dd\Phi_n(x,\theta,\bar\theta)\right]
    \exp\bigl(\ii I[\Phi]\bigr)
    \Phi_{n_1}(x_1,\theta_1,\bar\theta_1)
    \Phi_{n_2}(x_2,\theta_2,\bar\theta_2)\cdots\,.
  \tag{30.1.1}\label{eq:30.1.1}
\end{align}
where \(I[\Phi]\) is the action
\begin{equation}
  I[\Phi]
  =\frac12\int \dd^4x
    \left[\sum_n\Phi_n^*(x,\theta,\bar\theta)
      \Phi_n(x,\theta,\bar\theta)\right]_D
   +2\operatorname{Re}\int \dd^4x\,[f(\Phi)]_{\mathcal{F}}\,.
  \tag{30.1.2}\label{eq:30.1.2}
\end{equation}
(The factor \(1/2\) is introduced in the first term, as in
Eq.~\eqref{eq:26.4.3}, to give the component fields of \(\Phi\) a conventional
normalization.) But we cannot simply read off the supergraph Feynman rules
from Eq.~\eqref{eq:30.1.1}, because the functional integral over the
superfields \(\Phi_n\) must be constrained to satisfy the left-chiral
condition \(\bar D_{\dot\alpha}\Phi_n=0\).

This is analogous to a similar problem in electrodynamics. As discussed in
Section 12.3, at energies below the electron mass the interactions of soft
photons with each other are described by an effective action of the form:
\[
  I[f]
  =-\frac14\int \dd^4x\left[
    f_{\mu\nu}f^{\mu\nu}
    +c_1\bigl(f_{\mu\nu}f^{\mu\nu}\bigr)^2
    +c_2\bigl(\epsilon_{\mu\nu\rho\sigma}
       f^{\mu\nu}f^{\rho\sigma}\bigr)^2
  \right].
\]
But we cannot read off the Feynman rules from this formula without taking
into account the fact that the path integral is constrained by the
homogeneous Maxwell equations
\[
  \partial_\mu f_{\nu\rho}
  +\partial_\nu f_{\rho\mu}
  +\partial_\rho f_{\mu\nu}=0\,.
\]
As everyone knows, we deal with this constraint by introducing a four-vector
potential \(A_\mu\), with
\(f_{\mu\nu}=\partial_\mu A_\nu-\partial_\nu A_\mu\), so that the constraint
is automatically satisfied, and integrate over \(A_\mu(x)\), not
\(f_{\mu\nu}(x)\).

In the same way, we may adopt a trick already used in Section 26.6 to derive
the field equations for superfields, and introduce a non-chiral potential
superfield \(S_n(x,\theta,\bar\theta)\), with
\begin{equation}
  \Phi_n=\bar D^2S_n\,,
  \tag{30.1.3}\label{eq:30.1.3}
\end{equation}
where
\begin{equation}
  \bar D^2
  \equiv\bar D_{\dot\alpha}\bar D^{\dot\alpha}\,,
  \tag{30.1.4}\label{eq:30.1.4}
\end{equation}
so that \(\Phi_n\) obeys \(\bar D_{\dot\alpha}\Phi_n=0\), as required
for left chirality. We may therefore replace Eq.~\eqref{eq:30.1.1} by
the path-integral formula
\begin{align}
  &\left\langle
    T\left\{\Phi_{n_1}(x_1,\theta_1,\bar\theta_1),
      \Phi_{n_2}(x_2,\theta_2,\bar\theta_2),\ldots\right\}
   \right\rangle \notag\\
  &\qquad={}
  \int\left[\prod_{n,x,\theta,\bar\theta}
      \dd S_n(x,\theta,\bar\theta)\right]
    \exp\bigl(\ii I[\bar D^2S]\bigr)
    \bar D^2S_{n_1}(x_1,\theta_1,\bar\theta_1)
    \bar D^2S_{n_2}(x_2,\theta_2,\bar\theta_2)\cdots\,.
  \tag{30.1.5}\label{eq:30.1.5}
\end{align}
In expressing the action \eqref{eq:30.1.2} in terms of \(S_n\), we may recall
that the \(D\)-term of a superderivative does not contribute to the action, so
we may shift the operator
\((\bar D^2)^\dagger=D^2\) acting on \(S_n^*\) in the first term of
Eq.~\eqref{eq:30.1.2} to act instead on \(S_n\), which gives
\begin{equation}
  I[\bar D^2S]
  =\frac12\int \dd^4x
    \left[\sum_n S_n^*D^2\bar D^2S_n\right]_D
   +2\operatorname{Re}\int \dd^4x\,
    [f(\bar D^2S)]_{\mathcal{F}}\,.
  \tag{30.1.6}\label{eq:30.1.6}
\end{equation}
The \(\bar D^2\) operator in any one of the factors \(\bar D^2S\) in
any term in \(f(\bar D^2S)\) may be taken to act on the whole term, since
\(\bar D_{\dot\alpha}\) gives zero when acting on any other factor of
\(\bar D^2S\) in that term. In this way, we may write
\begin{equation}
  f(\bar D^2S)=\bar D^2\widetilde f(S)\,,
  \tag{30.1.7}\label{eq:30.1.7}
\end{equation}
where \(\widetilde f(S)\) is obtained from \(f(\bar D^2S)\) by omitting
any one operator \(\bar D^2\) in each term. For instance, for a single
type of superfield, if \(f(\Phi)=\sum_r c_r\Phi^r\), then
\[
  \widetilde f(S)=\sum_r c_rS(\bar D^2S)^{r-1}.
\]

Using Eqs.~\eqref{eq:30.1.6}, \eqref{eq:30.1.7}, and
\eqref{eq:26.3.31}, we may put the whole action in the form of a \(D\)-term
% With the ordered two-component derivatives,
% [\bar D^2 h]_{\mathcal F}=-2[h]_D.  The convention
% \Phi=\bar D^2S amounts to reversing the sign of the source potential
% superfield, so the interaction signs in (30.1.8)--(30.1.9) reverse
% while the quadratic term is unchanged.  Explicitly, for
% f=c(\bar D^2S)^r=\bar D^2[cS(\bar D^2S)^{r-1}], the superpotential
% contribution is 2 Re[f]_{\mathcal F}=-4 Re[\widetilde f]_D.
\begin{equation}
  I[\bar D^2S]
  =\frac12\int \dd^4x\sum_n
    [S_n^*D^2\bar D^2S_n]_D
   -4\operatorname{Re}\int \dd^4x\,[\widetilde f(S)]_D\,.
  \tag{30.1.8}\label{eq:30.1.8}
\end{equation}
Eq.~\eqref{eq:26.6.5} shows that this may also be written as a superspace
integral:
\begin{equation}
  I[\bar D^2S]
  =-\frac14\int \dd^4x\int \dd^4\theta
    \sum_nS_n^*D^2\bar D^2S_n
   +2\operatorname{Re}\int \dd^4x\int \dd^4\theta\,\widetilde f(S).
  \tag{30.1.9}\label{eq:30.1.9}
\end{equation}
This is the action we will use in the path-integral derivation of the
supergraph formalism.
